\scanpage{208}
SYSTÈMES de RACINES \mspage{(3)}\hfill Tribu n\textsuperscript{o} 59

linéaire réelle $\alpha\in\underline{\mathrm a}^{*}$ sur $\underline{\mathrm a}$, soit $\underline{\mathrm g}^{\alpha}=\{\mathrm x\in\underline{\mathrm g}\mid[\mathrm a,\mathrm x]=(\mathrm a).\mathrm x\;(\mathrm a\in\underline{\mathrm a})\}$. Alors l'ensemble $\mathrm R$ des éléments $\alpha\in\underline{\mathrm a}^{*}-0$ pour lesquels $\underline{\mathrm g}^{\alpha}\ne(0)$ est un système de racines, en général non réduit. Le groupe de Weyl de $\mathrm R$ s'identifie au groupe des automorphismes de $\underline{\mathrm a}$ induits par des éléments de $\mathrm{Ad}_{\underline{\mathrm g}}\,\underline{\mathrm k}$ laissant $\underline{\mathrm a}$ stable.

III. Soient $\mathrm k$ un corps, $\mathrm G$ un groupe linéaire algébrique connexe semi-simple sur $\mathrm k$, $\mathrm T$ un tore algébrique sur $\mathrm G$ décomposé sur $\mathrm k$ maximal, et $\mathrm X^{*}$ le groupe des caractères rationnels de $\mathrm T$. Le dual $\mathrm X_{*}$ de $\mathrm X^{*}$ s'identifie canoniquement au groupe des homomorphismes rationnels du groupe multiplicatif $\mathrm G_{\mathrm m}$ dans $\mathrm T$. Soit $\underline{\mathrm g}$ l'algèbre de Lie de $\mathrm G$ et, pour $\alpha\in\mathrm X^{*}$, soit $\underline{\mathrm g}^{\alpha}=\{\mathrm x\in\underline{\mathrm g}\mid\mathrm{Ad}\,\mathrm t(\mathrm x)=\alpha(\mathrm t).\mathrm x\;(\mathrm t\in\mathrm T)\}$. Alors l'ensemble $\mathrm R$ des éléments non nuls $\alpha\in\mathrm X^{*}$ pour lesquels $\mathrm g^{\alpha}\ne(0)$ est un système de racines, qui est réduit si $\mathrm T$ est aussi un sous-groupe de Cartan de $\mathrm G$, ce qui a en particulier toujours lieu lorsque $\mathrm k$ est algébriquement clos. Le groupe de Weyl $\mathrm W(\mathrm R)$ s'identifie au quotient du normalisateur de $\mathrm T$ par le centralisateur de $\mathrm T$ dans $\mathrm G$.

Si $\mathrm k=\underline{\underline{\mathrm C}}$ l'algèbre de Lie $\underline{\mathrm t}$ de $\mathrm T$ est une sous-algèbre de Cartan de $\underline{\mathrm g}$ et $\mathrm R$ est isomorphe au système de racines défini en I. Si $\mathrm k=\underline{\underline{\mathrm R}}$, alors l'ensemble des points réels de $\underline{\mathrm t}$ est conjugué à l'algèbre $\underline{\mathrm a}$ considérée en II, et $\mathrm R$ est isomorphe au système de racines défini en II.

IV. Soient $\mathrm G$ un groupe de Lie compact connexe semi-simple, $\underline{\mathrm g}$ son algèbre de Lie, $\mathrm T$ un tore maximal de $\mathrm G$ et $\mathrm X$ le groupe des homomorphismes continus de $\mathrm T$ dans $\underline{\underline{\mathrm C}}^{*}$. Pour tout $\alpha\in\mathrm X$, soit $\underline{\mathrm g}^{\alpha}=\{\mathrm x\in\underline{\mathrm g}\mid\mathrm{Ad}\,\mathrm t(\mathrm x)=\alpha(\mathrm t).\mathrm x\;(\mathrm t\in\mathrm T)\}$. Alors l'ensemble $\mathrm R$ des éléments $\alpha\in\mathrm X-0$ pour lesquels $\underline{\mathrm g}^{\alpha}\ne(0)$ est un système de racines réduit dans $\mathrm V=\mathrm X\otimes\underline{\underline{\mathrm Q}}$, dont le groupe de Weyl s'identifie au quotient par $\mathrm T$ du normalisateur de $\mathrm T$ dans $\mathrm G$. On a $\mathrm Q(\mathrm R)\subset\mathrm X\subset\mathrm P(\mathrm R)$ et $\mathrm X=\mathrm P(\mathrm R)$, (resp. $\mathrm X=\mathrm Q(\mathrm R)$), si $\mathrm G$ est simplement connexe, (resp. isomorphe à $\mathrm{Ad}\,\underline{\mathrm g}$). L'algèbre de
% archive sheet number: 207

